\section{An estimate on the shrinking spheres}
\label{sec:angular}

The functional in Proposition~\ref{balls:functional} is defined at every positive
radius.  We now estimate its representing density uniformly as the radius
shrinks.  The equation for this density has second-order angular and base
terms.  Comparing a degree-raising operator with a degree-lowering operator
will control both kinds of derivatives; a second comparison will retain the
principal terms caused by the changing metric.

\subsection{The equation for the density}

We use the reference metric $\gamma$, the normalized second principal
matrix $Q$, and $h=Q-\gamma^{-1}$ from
\eqref{balls:normalization}.  In this section the subscript $\delta$ on the positive
radius $r$ is suppressed.  Put
\[
 f=\log r,\qquad \beta=|df|_\gamma.
\]
The required profile bounds are
\begin{equation}\label{angular:profile-bounds}
 \beta\ge\beta_0,\qquad
 \operatorname{Hess}_\gamma f\le-c_0\beta\gamma,
 \qquad |\nabla^j f|\le C_j\beta^j\quad(1\le j\le8),
\end{equation}
where $c_0>0$ and $C_j$ are fixed.  Lemma~\ref{balls:profiles} provides
these bounds with $\beta_0$ arbitrarily large by choosing $p$
sufficiently large.  All background metric bounds
below are on a fixed relatively compact coordinate region; their constants
are uniform in the spatial dilation.

Let $S_\gamma Y$ be the bundle of $\gamma$-unit tangent spheres over the
base region $Y$.  Each sphere has probability measure $d\sigma_y$.
Horizontal differentiation $\nabla$ uses parallel transport for $\gamma$,
including the covariant tensor indices of an operator or section.  This
connection preserves the sphere measure.  Unless otherwise specified, all
norms and full adjoints use $dV_\gamma(y)\,d\sigma_y(\omega)$.

On each tangent unit ball, extend a sphere function harmonically for the
Euclidean metric $\gamma(y)$.  Denote the boundary radial derivative of
this extension by $B$, and its ambient derivatives in an orthonormal frame
by $A_i$.  Thus $B=k$ on the space $\mathcal H_k$ of spherical harmonics
of degree $k$, and $A_i:\mathcal H_k\to\mathcal H_{k-1}$.  The star in
$A_i^*$ denotes the sphere adjoint.  These operators are parallel, with
the indicated tensor indices, and
\begin{equation}\label{angular:elementary-identities}
 [B,A_i^*]=A_i^*,\qquad \sum_i A_i^*A_i^*=0.
\end{equation}
Indeed differentiation lowers homogeneous degree by one, and the second
identity is the adjoint of the vanishing Laplacian of a harmonic extension.
For $s\in\mathbb R$, use the fiber Sobolev norm
\[
 \|v(y,\cdot)\|_s=\|(1+B)^s v(y,\cdot)\|_{L^2(S_y)}.
\]
It is equivalent to the usual Sobolev norm on the unit sphere.  Tensor-valued
versions use the sum over an orthonormal frame.

We write $\Psi^m$ for the classical pseudodifferential operators of order
$m$ on the sphere.  For a positive function $a(y)$, the notation
$a\Psi^m$ means that, after division by $a$, the finitely many symbol and
operator seminorms used below are bounded.  A statement with $j$ base
derivatives permits the indicated additional factor $\beta^j$.

\begin{lemma}[Moving trace operators]\label{angular:trace}
Suppose that the normal balls and the functional $T_y$ of
Proposition~\ref{balls:functional} are defined, and that
\begin{equation}\label{angular:smallness}
 \sum_{|\alpha|\le8}|\partial_y^\alpha h|\le\eps r,
 \qquad r\beta^2\le\eps.
\end{equation}
If $u$ is first-metric harmonic near one such ball, set
\[
 U(y,\omega)=u\bigl(\exp_y^\gamma(r(y)\omega)\bigr).
\]
There are sphere operators $\mathcal C_i$ such that
$\nabla_iU=\mathcal C_iU$.  Their sphere adjoints satisfy
\begin{equation}\label{angular:C-expansion}
 \mathcal C_i^*=r^{-1}A_i^*+f_iB+E_i+F_i,
 \qquad E_i\in r\Psi^1,\quad F_i\in\Psi^0,
 \qquad f_i=\nabla_i f.
\end{equation}
The normalized remainders $r^{-1}E_i$ and $F_i$ have uniform symbol
bounds after up to six base derivatives, with the allowance $C_j\beta^j$.

Write the normalized second equation, in the reference connection, as
$Q^{ij}\nabla_i\nabla_j+b_2^i\nabla_i$.  Its drift $b_2$ is uniformly
controlled.  If $\psi(y,\omega)$ represents $T_y$ against $d\sigma_y$,
then
\begin{equation}\label{angular:P-equation}
 P\psi=0,\qquad
 P=r\left[Q^{ij}(\nabla_i+\mathcal C_i^*)
                    (\nabla_j+\mathcal C_j^*)
             +b_2^i(\nabla_i+\mathcal C_i^*)\right].
\end{equation}
The products in this formula are covariant products.  In particular,
$Q$ is not differentiated by taking a base adjoint.
\end{lemma}

\begin{proof}
Normalize the principal part of the first equation to $\gamma^{-1}$ and
pull it back by $z\mapsto\exp_y^\gamma(rz)$ to the unit ball.  After
multiplication by $r^2$, its principal coefficients are
$\delta^{ij}+O(r^2)$ and its first-order coefficients are $O(r)$.
The resulting smooth family remains uniformly elliptic and has an
invertible Dirichlet problem, including at $r=0$.  Its radial derivative
operator therefore has the expansion
\begin{equation}\label{angular:radial-expansion}
 B+r B_1(y)+r^2 B_2(y,r),\qquad
 B_1\in\Psi^0,\quad B_2\in\Psi^1.
\end{equation}
Here is the parameter justification.  The classical boundary-symbol
construction is described in
\citep[Theorem~1.1(i) and Proposition~3.1]{JoshiLionheart2005};
we keep the center and radius dependence explicitly.
In boundary coordinates the decaying
normal root of the principal symbol constructs the Dirichlet parametrix.
Successive lower-order transport equations construct its symbols smoothly
in $(y,r)$.  The first normal derivative trace is an operator of order one,
whose principal symbol depends only on the principal coefficients and the
differentiating vector.  Those coefficients have zero first $r$-derivative
at $r=0$, so that derivative of the trace operator has order zero.  Taylor's
formula in this symbol class gives \eqref{angular:radial-expansion} to
smoothing error.  The true solution differs from the parametrix by the
Dirichlet inverse of a smooth error; differentiated Dirichlet estimates
control this remainder in every fixed seminorm.  More explicitly, on a
compact center set the rescaled Dirichlet operator is smoothly
invertible between each fixed pair of Sobolev spaces, uniformly for
$0\le r\le r_0$.  Differentiating that inverse controls every fixed
number of parameter derivatives of the smoothing correction; increasing
the Sobolev orders controls its smooth kernel seminorms.  Thus the
Taylor expansion holds in the complete classical symbol topology,
including the smoothing remainder, and not merely in one Sobolev
operator norm.  This proves the stated expansion and its parameter bounds.

Differentiate the exponential map in its center, transporting $\omega$
parallel to the center velocity.  The Jacobi field has initial derivative
zero.  After inverting the differential in the $z$ variable, the velocity
acting on the unit-ball solution is
\[
 r^{-1}\bigl(e_i+O(r^2)\bigr)+f_i\omega.
\]
The coefficient error is smooth in $(y,r,\omega)$.  Tangential derivatives
act directly on the boundary value; the radial derivative is
\eqref{angular:radial-expansion}.  At $r=0$ the first term is $r^{-1}A_i$.
The order-one errors are $O(r)$, and the order-zero errors are bounded:
in particular $r f_i B_1$ is bounded because $r\beta\le\eps/\beta$.
Taking sphere adjoints gives \eqref{angular:C-expansion}.  Differentiating
the radius costs at most one factor $\beta$ per derivative, by
\eqref{angular:profile-bounds}; the asserted parameter estimates follow.

For fixed $u$, differentiate
$T_yu=\langle\psi(y,\cdot),U(y,\cdot)\rangle_{S_y}$.
Metric compatibility and $\nabla_iU=\mathcal C_iU$ put the differentiated
operator on $\psi$ as $\nabla_i+\mathcal C_i^*$.
The local second-metric harmonicity in
Proposition~\ref{balls:functional} proves that
\eqref{angular:P-equation} pairs to zero with the trace of every
function harmonic near the fixed closed ball.  No global extension
or source-family approximation is used in this step.
At a fixed sphere these traces are dense among
smooth boundary data: solve the smooth Dirichlet problem on slightly
enlarged balls, using boundary values transported from the limiting sphere,
and use smooth dependence of the solution on the radius.  Thus the paired
equation is the stated equation for $\psi$.
\end{proof}

\subsection{The model operator and its principal perturbations}

We separate a degree-raising model operator from the remaining
principal and lower-order terms of the moving trace equation.  We continue under the hypotheses of
Lemma~\ref{angular:trace}, including \eqref{angular:smallness}.

Define
\begin{equation}\label{angular:T-definition}
 D_i^+=\nabla_i+f_iB,\qquad D_i^-=-\nabla_i+f_iB,
 \qquad
 T=\sum_i A_i^*D_i^+,\quad T^\flat=\sum_i A_iD_i^-.
\end{equation}
The symbol $T$ here denotes a differential operator on sphere functions;
the transfer functional continues to be denoted $T_y$.
Expand \eqref{angular:P-equation} by
\eqref{angular:C-expansion}.  The identities
\eqref{angular:elementary-identities} give
\[
 D_i^+(r^{-1}A_j^*)=r^{-1}A_j^*D_i^+,
 \qquad \sum_iA_i^*A_i^*=0.
\]
Thus, in ordinary base coordinates and an orthonormal sphere
trivialization,
\begin{equation}\label{angular:P-decomposition}
 P=2T+U+J,
\end{equation}
where $J$ is a uniformly bounded family in $\Psi^1$, with no base
derivatives, and
\begin{equation}\label{angular:U-bounds}
 U=\sum_{|\alpha|\le2}U_\alpha(y)\partial_y^\alpha,
 \qquad U_\alpha\in\eps\beta^{-|\alpha|}
                               \Psi^{2-|\alpha|}.
\end{equation}
Up to five derivatives of these coefficients have the additional
allowance $C_j\beta^j$.  The second-base-derivative term is precisely
\begin{equation}\label{angular:real-principal-part}
 rQ^{ij}\partial_{y_i}\partial_{y_j};
\end{equation}
its coefficients are real scalars independent of the sphere variable.

We detail these bounds.  The pure angular second derivative contracted
with $\gamma^{-1}$ vanishes.  The remaining contraction is
$(h^{ij}/r)A_i^*A_j^*$, and belongs to the zero-base-derivative term of
\eqref{angular:U-bounds}.  The mixed correction is
$2h^{ij}A_i^*D_j^+$ and has the stated small bounds.  The $D^+D^+$ term
has coefficient $r$; its base orders two, one and zero have coefficients
$O(r)$, $O(r\beta)$ and $O(r\beta^2)$, respectively.  The condition
$r\beta^2\le\eps$ gives \eqref{angular:U-bounds} for all three.
Terms containing $E_i$ obey the same bounds.  Terms containing $F_i$
also do so, except for their cross terms with $r^{-1}A_j^*$, which are
bounded angular order-one terms and belong to $J$.  The drift term with
$A_i^*$ belongs to $J$ as well.  Differentiating any of these expressions
gives the stated coefficient bounds by
\eqref{angular:profile-bounds} and \eqref{angular:smallness}.
A horizontal connection written in this trivialization adds a bounded
angular first derivative, which preserves the same estimates.  The bounds
for $J$ use the fixed background and remain uniform as $\beta_0$ is
increased.

The decomposition identifies the model operator $2T$ and the
second-order perturbation $U$.  To carry the model estimate over to $P$
we shall need more than an ordinary first-derivative estimate: the error forms
also contain mixed base and angular derivatives.  For a sphere function
$\varphi$, put
\begin{equation}\label{angular:N-definition}
 \mathcal N(\varphi)^2
 =\int_Y\left(
  \|\nabla\varphi\|_0^2+\beta^2\|\varphi\|_1^2
  +\beta^{-1}\|\nabla\varphi\|_{1/2}^2
  +\beta\|\varphi\|_{3/2}^2\right)dV_\gamma.
\end{equation}

The last two terms of $\mathcal N$ place half an angular derivative
on $\nabla\varphi$ and three halves on $\varphi$.  They will bound the
remaining error forms left by the comparison with $P$.

\subsection{A comparison of raising and lowering operators}

The next estimate supplies these two derivative strengths for the model
operator.  Spherical harmonics are trace-free symmetric tensors, and
comparing their raising and lowering operators is familiar from energy
identities in tensor tomography; see, for example, \cite{PSU15,GPSU16}.
Here the negative Hessian of the spatial weight supplies the positive
term that absorbs the bounded curvature.  We give the calculation with
its degree factors.

\begin{lemma}[Weighted raising estimate]\label{angular:raising}
Let $n\ge3$, and suppose \eqref{angular:profile-bounds} holds on a fixed
base region with uniformly bounded reference geometry.  There are
$\beta_0$ and $c>0$, depending only on that geometry and the constants in
\eqref{angular:profile-bounds}, such that every smooth, compactly
base-supported function with zero fiber mean satisfies
\begin{equation}\label{angular:raising-estimate}
 \|T\varphi\|^2-\|T^\flat\varphi\|^2
       \ge c\,\mathcal N(\varphi)^2.
\end{equation}
\end{lemma}

\begin{proof}
Both operators change angular degree by exactly one, so it suffices to
prove the estimate at one input degree $k\ge1$.  Identify that component
with a harmonic homogeneous polynomial $v(z)$ of degree $k$.  Use the
Fischer inner product, in which the monomials $z^\alpha$ are orthogonal
with squared norm $\alpha!$.  Polynomial differentiation $a_i=\partial_{z_i}$
is adjoint to multiplication $M_i=z_i$.  On harmonic polynomials of degree
$k$ the Fischer squared norm is the normalized squared sphere norm
multiplied by
\[
 N_k=n(n+2)\cdots(n+2k-2),\qquad N_0=1.
\]
For completeness, integration against a standard real Gaussian identifies
the Fischer norm on trace-free homogeneous polynomials with their
Gaussian $L^2$ norm;
integrating radially then gives the factor $N_k$.  These classical
identities are also recorded in \cite[Equation~(2.1) and
Theorem~2.1]{Render08}.

Set $d=k+n/2-1$ and use $D_i^\pm=\pm\nabla_i+k f_i$ at this fixed degree.
For harmonic polynomials $q_i$ of degree $k$, the Laplacian of
$\sum_iM_iq_i$ is $2\sum_i a_iq_i$.  Since
\[
 \Delta_z(|z|^2s)=4d\,s\qquad(s\in\mathcal H_{k-1}),
\]
its harmonic projection is
$\sum_iM_iq_i-|z|^2\sum_i a_iq_i/(2d)$.  Orthogonality of this projection
and its complement shows that its squared norm is
\[
 \sum_i\|q_i\|^2+\sum_{i,j}(a_jq_i,a_iq_j)
             -d^{-1}\Bigl\|\sum_i a_iq_i\Bigr\|^2.
\]
Changing from Fischer adjoints to sphere adjoints gives
\[
 A_i^*|_{\mathcal H_k}
   =\frac{N_{k+1}}{N_k}\Pi_{k+1}M_i,
\]
where $\Pi_{k+1}$ is harmonic projection in homogeneous degree $k+1$.
Consequently $N_k$ times the left side of
\eqref{angular:raising-estimate} is
\begin{equation}\label{angular:Fock-difference}
 \begin{split}
 (2d+2)\left[\|D^+v\|^2
   +\sum_{i,j}(a_jD_i^+v,a_iD_j^+v)
   -\frac1d\Bigl\|\sum_i a_iD_i^+v\Bigr\|^2\right]
 -2d\Bigl\|\sum_i a_iD_i^-v\Bigr\|^2.
 \end{split}
\end{equation}
All norms in this calculation include integration over the base and use
the polynomial inner product in the fibers.  The factors are
$N_{k+1}/N_k=2d+2$ and $N_k/N_{k-1}=2d$.

Write $C_{ij}=2k(\operatorname{Hess}_\gamma f)_{ij}$ and
\[
 C[av]=\int_Y\sum_{i,j}C_{ij}\langle a_iv,a_jv\rangle\,dV_\gamma.
\]
Covariant integrations by parts give
\begin{align}
 \sum_{i,j}(a_jD_i^+v,a_iD_j^+v)
   &=\Bigl\|\sum_i a_iD_i^-v\Bigr\|^2-C[av]
                                    +O(k^2)\|v\|^2,\label{angular:ibp-one}\\
 \Bigl\|\sum_i a_iD_i^+v\Bigr\|^2
   &\le k\|D^-v\|^2-C[av]+O(k^2)\|v\|^2,\label{angular:ibp-two}\\
 \|D^-v\|^2
   &=\|D^+v\|^2+\int_Y\tr C\,|v|^2\,dV_\gamma.
                                                    \label{angular:ibp-three}
\end{align}
Here and below a form denoted $O(k^2)\|v\|^2$ has absolute value at most
a fixed constant times that quantity.  To verify the first two formulas,
move a $D_j^+$ across the pairing and commute
$D_j^-D_i^+$ to $D_i^+D_j^--C_{ij}$.  The reordered term in
\eqref{angular:ibp-two} is at most $k\|D^-v\|^2$, because
$\sum_i|a_iw|^2=k|w|^2$ for $w\in\mathcal H_k$.
The extra commutator is curvature: it acts as a sum of rotations of norm
$O(k)$ at degree $k$, and the two contractions supply one more factor $k$.
This proves the claimed error bounds.  Expanding the two squares proves
\eqref{angular:ibp-three}.  For complex functions all the displayed forms
are understood through their real parts.

Substituting \eqref{angular:ibp-one} into
\eqref{angular:Fock-difference} rewrites that expression as
\begin{equation}\label{angular:after-first-ibp}
 \begin{split}
 &(2+2/d)\left[d\|D^+v\|^2
       -\Bigl\|\sum_i a_iD_i^+v\Bigr\|^2-dC[av]\right]\\
 &\qquad+2\Bigl\|\sum_i a_iD_i^-v\Bigr\|^2
                      +O(k^3)\|v\|^2.
 \end{split}
\end{equation}
We must control the remaining negative divergence square while retaining
both an ordinary derivative estimate and an additional half-order
angular derivative estimate.  The latter requires a coefficient of
size $k/\beta$ in front of $|D^+v|^2$, including when $k$ is much larger
than $\beta$.  Choose a small constant $\vartheta>0$.
On the fraction $1-\vartheta$ of its negative divergence square use
\eqref{angular:ibp-two}--\eqref{angular:ibp-three}.  On the remaining
fraction use the pointwise convex combination
\begin{equation}\label{angular:convex-bound}
 \begin{split}
 \Bigl|\sum_i a_iD_i^+v\Bigr|^2
 &\le(1-\zeta)(k+n-1)|D^+v|^2\\
 &\quad+\zeta\left(2\Bigl|\sum_i a_iD_i^-v\Bigr|^2
                         +8k^3\beta^2|v|^2\right),
 \qquad \zeta=c_1/\beta.
 \end{split}
\end{equation}
The fraction $1-\vartheta$ in the integrated estimates is constant;
the $y$-dependent $\zeta$ is used only in this pointwise inequality.
The first bound follows from the adjoint multiplication operators, since
the contraction $(q_i)\mapsto\sum_i a_iq_i$ has squared norm at most
$k+n-1$ on polynomials of degree $k$.  The second follows from
$D_i^++D_i^-=2k f_i$ and
$|\sum_i f_i a_i v|^2\le k\beta^2|v|^2$.

Combining these bounds in \eqref{angular:after-first-ibp}, the coefficient
retained for $|D^+v|^2$ is
\begin{equation}\label{angular:positive-degree-factor}
 (2+2/d)\bigl[d-k-\vartheta(n-1)
                         +\vartheta\zeta(k+n-1)\bigr].
\end{equation}
Since $d-k=(n-2)/2$, a sufficiently small fixed $\vartheta$ makes this
at least a fixed positive constant plus a positive multiple of
$c_1k/\beta$.  The coefficient of
$|\sum_i a_iD_i^-v|^2$ is nonnegative when $\beta_0$ is sufficiently
large.  The Hessian contribution is exactly
\begin{equation}\label{angular:hessian-contribution}
 -(2+2/d)\left[(d-1+\vartheta)C[av]
       +(1-\vartheta)k\int_Y\tr C\,|v|^2\,dV_\gamma\right].
\end{equation}
For $k\ge1$ and $n\ge3$, both coefficients inside this expression are
positive.  The Hessian hypothesis bounds it below by
$c k^3\int\beta|v|^2$.  This absorbs the curvature error $O(k^3)\|v\|^2$
by increasing $\beta_0$, and absorbs the last term of
\eqref{angular:convex-bound} by choosing $c_1$ sufficiently small.
We have therefore retained positive multiples of
\[
 \|D^+v\|^2,\qquad
 \int_Y\frac{k}{\beta}|D^+v|^2\,dV_\gamma,
 \qquad k^3\int_Y\beta|v|^2\,dV_\gamma.
\]
The first controls both $\|\nabla v\|^2$ and
$k^2\int\beta^2|v|^2$, since its cross term is
$-k\int\Delta_\gamma f\,|v|^2\ge0$.  For the weighted derivative use the
pointwise inequality
\[
 \frac{k}{\beta}|\nabla v|^2
 \le2\frac{k}{\beta}|D^+v|^2+2k^3\beta|v|^2.
\]
No differentiation of $\beta^{-1}$ is required.  Dividing by $N_k$, and
then summing the orthogonal degree components, proves
\eqref{angular:raising-estimate}; the weights $(1+k)^s$ are comparable
to $k^s$ because the zero degree is excluded.
\end{proof}

\subsection{Keeping the principal perturbations}

The raising estimate controls the model operator.  The actual equation
\eqref{angular:P-equation} also contains second base derivatives and a
small second-order angular perturbation.  We retain these terms in a
comparison of squared norms.

\begin{proposition}[Coercivity for the moving trace equation]
\label{angular:coercivity}
Fix a base region with bounded reference geometry and the constants in
\eqref{angular:profile-bounds}.  Suppose $P$ is the operator
\eqref{angular:P-equation}, with the trace operators of
Lemma~\ref{angular:trace}.  There exist $\beta_0$ sufficiently large,
$\eps_0>0$ and $C<\infty$, depending only on these fixed bounds, such that
\eqref{angular:smallness} with $\eps\le\eps_0$ implies
\begin{equation}\label{angular:coercive-estimate}
 \mathcal N(\varphi)\le C\|P\varphi\|
\end{equation}
for every smooth, compactly base-supported function $\varphi$ with
zero mean on each sphere.  These constants are independent of the
shrinking parameter.  Only the eight indicated derivatives of the small
tensor $h$ are required.
\end{proposition}

\begin{proof}
Let a dagger denote the full adjoint for $dV_\gamma d\sigma$.
Covariant integration and the degree shift give
\begin{equation}\label{angular:Z-definition}
 Z:=T^\flat-T^\dagger=\sum_i f_iA_i.
\end{equation}
A direct expansion yields
\begin{equation}\label{angular:norm-comparison}
 \begin{split}
 &\|(2T+U)\varphi\|^2-
                \|(2T^\flat+U^\dagger)\varphi\|^2\\
 &\qquad=4\bigl(\|T\varphi\|^2-\|T^\flat\varphi\|^2\bigr)
                       +\langle\mathcal E\varphi,\varphi\rangle,
 \end{split}
\end{equation}
where
\begin{equation}\label{angular:error-operator}
 \mathcal E=2[T^\dagger,U]+2[U^\dagger,T]+[U^\dagger,U]
                       -2\bigl(UZ+(UZ)^\dagger\bigr).
\end{equation}
We will prove
\begin{equation}\label{angular:error-form-bound}
 |\langle\mathcal E\varphi,\varphi\rangle|
                         \le C\eps\mathcal N(\varphi)^2.
\end{equation}

Here are the complete order and weight counts needed for this claim.
An operator has count $(m,s)$ if its coefficient at $\partial_y^\alpha$
is angular order $m-|\alpha|$, of size $O(\beta^{s-|\alpha|})$,
with the derivative allowances already specified.
Adjoints preserve the count, and products add the counts.  When a base
derivative hits a coefficient it consumes one differential order and
adds a factor $\beta$, so it preserves the weight index.  A commutator
of scalar angular operators loses one angular order: the products of
their principal symbols cancel.  Hence a scalar commutator loses one in
total order while preserving the sum of the weight indices.
The counts of $T,T^\dagger,U,Z$ are, respectively,
\[
 (2,1),\quad(2,1),\quad(2,0),\quad(1,1),
\]
and every occurrence of $U$ carries the factor $\eps$.  Thus
\eqref{angular:error-operator} has count at most $(3,1)$ and a factor
$O(\eps)$; the quadratic $U$ term is smaller.

There is a further cancellation that the total count alone does not
express: no third base derivative survives.  To check it explicitly,
use an orthonormal sphere trivialization and write the base measure as
$\mu(y)\,dy$.  The sphere probability measure is then fixed.  Write
\[
 U=a^{ij}\partial_i\partial_j+\mathcal B^i\partial_i+\mathcal C,
 \qquad a^{ij}=rQ^{ij}=a^{ji}\in\mathbb R.
\]
The coefficients $\mathcal B^i$ and $\mathcal C$ are angular
operators.  Replacing a horizontal derivative by a coordinate derivative
plus its angular connection term changes only base orders at most one,
so these are all the second-base coefficients.  If a star denotes the
sphere adjoint, direct integration by parts gives
\begin{align*}
 U^\dagger-U
 ={}&\left(2\mu^{-1}\partial_j(\mu a^{ij})
                -\mathcal B^i-(\mathcal B^i)^*\right)\partial_i\\
 &+\mu^{-1}\partial_i\partial_j(\mu a^{ij})
       -\mu^{-1}\partial_i\bigl(\mu(\mathcal B^i)^*\bigr)
       +\mathcal C^*-\mathcal C.
\end{align*}
In particular it has base order at most one.
In $[T^\dagger,U]$ and
its partner, the possible third-base-derivative coefficient is a commutator of an angular coefficient
with $rQ^{ij}$, which is zero by
\eqref{angular:real-principal-part}.  For $[U^\dagger,U]$, first write
it as $[U^\dagger-U,U]$.  The operator $U^\dagger-U$ has no second base
derivative, since the coefficient in
\eqref{angular:real-principal-part} is real.  Taking adjoints against
the smooth base density only produces lower base orders.  Its remaining
possible third-base-derivative coefficients commute for the same reason.
Finally, $UZ$ already has at most two base derivatives.
It follows that $\mathcal E$ is a sum of terms of precisely the following
three permitted types:
\begin{equation}\label{angular:error-types}
 \eps\beta^{-1}\Psi^1\partial_y^2,
 \qquad \eps\Psi^2\partial_y,
 \qquad \eps\beta\Psi^3.
\end{equation}
This argument also covers lower orders, since $\beta\ge1$.

For the first type in \eqref{angular:error-types}, integrate once in the
base.  The principal form is bounded by
$C\eps\int\beta^{-1}\|\partial_y\varphi\|_{1/2}^2$.
A derivative of its coefficient costs one factor $\beta$ and leaves an
$\eps\Psi^1\partial_y$ term.  This includes differentiation of the
weight itself: $|d\beta|\le C\beta^2$ implies
$|d(\beta^{-1})|\le C$.  The second type has the form bound
\[
 C\eps\int_Y
     \|\partial_y\varphi\|_{1/2}\|\varphi\|_{3/2}\,dV_\gamma
 \le C\eps\int_Y\left(
     \beta^{-1}\|\partial_y\varphi\|_{1/2}^2
                 +\beta\|\varphi\|_{3/2}^2\right)dV_\gamma.
\]
The third type is bounded by
$C\eps\int\beta\|\varphi\|_{3/2}^2$.
The lower-order term from differentiating a coefficient satisfies the
same bound.  Replacing coordinate derivatives by horizontal derivatives
adds a bounded angular first derivative, again controlled by these norms.
A fixed finite partition of the base region has the same harmless
lower-order errors.  This proves \eqref{angular:error-form-bound}.

The constants in the preceding estimates are uniform for all
$\beta\ge\beta_0$ once a fixed lower threshold is met.  Choose
$\eps_0$ small relative to the constant in
Lemma~\ref{angular:raising}.  Dropping the nonnegative second squared
norm in \eqref{angular:norm-comparison} gives
$\mathcal N(\varphi)\le C\|(2T+U)\varphi\|$.
Finally,
\[
 \|J\varphi\|\le C\|\varphi\|_{L^2(Y;H^1(S_y))}
                         \le C\beta_0^{-1}\mathcal N(\varphi),
\]
so a sufficiently large $\beta_0$ absorbs $J$ and proves
\eqref{angular:coercive-estimate}.

All these operations use angular symbol estimates pointwise in the base,
and ordinary base differential operators of degree at most two.  Formal
adjoints, the commutator expansion and the one integration by parts use
at most the five base coefficient derivatives specified in
\eqref{angular:U-bounds}.  Construction of these coefficients differentiates
the moving trace operators once and the second metric to only finitely
many orders below eight.  Higher angular symbol seminorms depend on the
fixed smooth first-metric family; differentiating in the angular variable
does not differentiate the center-dependent tensor $h$.  Thus the eight
smallness derivatives in \eqref{angular:smallness} suffice in every fixed
dimension.  No smallness assumption on all derivatives of $h$ is used.
\end{proof}

\subsection{The constant mode and the uniform density bound}

We now apply the coercive estimate to the transfer density constructed in
Section~\ref{sec:balls}, restoring the notation $r=r_\delta$.
Its normalization is $\int_{S_y}\psi\,d\sigma_y=1$, so
$\chi(\psi-1)$ has zero fiber mean for the fixed cutoff
\eqref{balls:cutoff-region}.  This is the mean-zero condition required
by Proposition~\ref{angular:coercivity}.

The decomposition \eqref{angular:P-decomposition} gives a uniform
bound for $P1$: its raising part annihilates $1$, only $U$'s bounded
zero-base-derivative coefficient acts on $1$, and $J$ is uniformly
bounded in $\Psi^1$.  There is no factor growing with $\beta$ in this
assertion.

For any fixed smooth base cutoff, the same decomposition gives
\begin{equation}\label{angular:cutoff-bound}
 \|[P,\chi]w\|
 \le C_\chi\left(
  \|w\|_{L^2(\supp d\chi;H^1(S_y))}
    +\|r\nabla w\|_{L^2(\supp d\chi\times S_y)}\right),
\end{equation}
where the support can be replaced by a fixed small neighborhood of the
supports of all derivatives of $\chi$.
Indeed $[2T,\chi]=2\sum_iA_i^*(\nabla_i\chi)$, the second base derivative
of $U$ contributes $2rQ^{ij}(\partial_i\chi)\partial_j$ and a bounded
zeroth-base-derivative term, and its first base derivative contributes a
bounded angular order-one term.  The operator $J$ commutes with $\chi$.

\begin{lemma}[Uniform transfer density]\label{local:density}
Fix the structural parameters so that Propositions
\ref{balls:continuation} and \ref{angular:coercivity} apply.  There is
a constant $C$ independent of $\delta$ and of sufficiently small
spatial dilations $\lambda$ such that
\begin{equation}\label{local:density-bound}
 \|\chi\psi\|_{L^2(S_\gamma Y)}\leq C
\end{equation}
whenever the coefficient condition
\eqref{balls:bootstrap-condition} holds.
\end{lemma}

\begin{proof}
Set $\varphi=\chi(\psi-1)$, which is compactly supported in the base
and has zero fiber mean.  For the actual moving-sphere operator $P$,
Equation~\eqref{angular:P-equation} gives
\begin{equation}\label{local:density-equation}
 P\varphi=P\bigl(\chi(\psi-1)\bigr)
       =-\chi P1+[P,\chi](\psi-1).
\end{equation}
The cutoff estimate \eqref{angular:cutoff-bound} involves only angular
derivatives through order one and $r\nabla\psi$ on the supports of
derivatives of $\chi$.  Proposition~\ref{balls:functional} bounds
these there uniformly: on the lower region the density is that of
ordinary evaluation, and on the other region the radii stay at a fixed
positive separation.  Together with the bound for $P1$, this gives a
uniformly bounded $L^2$ norm for the right side of
\eqref{local:density-equation}.
Proposition~\ref{angular:coercivity} bounds $\varphi$; adding the fixed
function $\chi$ proves \eqref{local:density-bound}.
\end{proof}

The uniform bound now applies to the density representing the actual
Green functional at every radius allowed by
\eqref{balls:bootstrap-condition}.  Exact affine moments will turn it
into a two-power bound for the paired harmonic differences.
